\section{第三天}
\begin{enumerate}
  \item 必读提示。
        中国特色社会主义法治道路直接关联了中国式法治现代化这个热到冒泡的热门话题。
        概言之，中国式现代化建设离不开法治建设，
        而法治建设要保持中国特色的最根本要求，
        就是坚持从中国实际出发，
        坚持中国特色社会主义法治道路。
        所以题目中一旦\bs{涉及中国式现代化或中国特色等主题词}，
        都是从这部分内容来获取答题资源的。
        这个知识点不背，
        不仅是行业之幸，
        甚至可以说是行业万幸！！！
        其中核心要义的第 2 点，
        即坚持中国特色社会主义制度的阐述主要集中在\bs{对社会主义制度优越性的描述}，不需要全都背，
        以能够阐述 3 - 5 句为目标即可（可以正反论述凑字数）。
        背！必须背！往死里背！
  \item 政治方向之\bs{坚持以人民为中心}（哪怕背成脑瘫，也要背下来）。
        \begin{enumerate}
          \item 概念和价值（3 个根本立场）。
                \begin{enumerate}
                  \item 是中特法的\bs{根本立场}。
                  \item 全治国\zx{最广泛、最深厚的基础}是人民，
                        推全国的\zx{根本目的}是依法保障人民权益。
                  \item 人民群众是我们党的\bs{力量源泉}，人民立场是中国共产党的\bs{根本政治立场}，
                  \item 也是\bs{新时代坚持和发展社会主义的根本立场}，
                  \item 是中特法的\bs{本质要求}。来个总结，首尾呼应。
                  \item 坚持以人民为中心，
                        \begin{enumerate}
                          \item 深刻回答了推全国，建法国\zx{为了谁、依靠谁}的问题。
                          \item \zx{始终代表}最广大\bs{人民}根本利益，保证人民\zx{当家}作主，体现人民共同\zx{意志}，维护人民合法\zx{权益}，
                          \item 是我国\zx{国家制度和国家治理体系}的\bs{本质属性}，
                                也是国家制度和国家治理体系\zx{有效运行、充满活力}的\bs{根本所在}。
                        \end{enumerate}
                \end{enumerate}
          \item 如何（左边关键句必背）。
                \begin{enumerate}
                  \item 坚持人民的\bs{主体地位}。
                        \begin{enumerate}
                          \item 确保人民群众\zx{享受当家作主}权利的同时，还要\zx{承担}应尽的义务和责任，
                          \item 当法治事业的\zx{忠实的崇尚者、自觉的追随者、坚定的捍卫者}。
                        \end{enumerate}
                  \item 牢牢把握\bs{社会公平正义}这一法治价值追求。
                        \begin{enumerate}
                          \item 把公平正义\zx{贯穿到立执司守}的全过程和各方面，
                          \item 努力让人民群众在每一项法律制度、每一个执法决定、每一宗司法案件中\zx{都感受到}公平正义。
                        \end{enumerate}
                  \item 推全国的\bs{根本目的}是依法保障人民权益。
                        \begin{enumerate}
                          \item 必须坚持法治\zx{为了}人民、\zx{依靠}人民、\zx{造福}人民、\zx{保护}人民，
                          \item 系统研究、谋划和解决法治领域人民群众\zx{反映强烈}的突出问题，
                          \item 不断增强人民群众\zx{获得感、幸福感、安全感}，\zx{用法治保障}人民\zx{安居乐业}。
                        \end{enumerate}
                \end{enumerate}
        \end{enumerate}
  \item 政治方向之\bs{坚持中国特色社会主义法治道路}。
        \begin{enumerate}
          \item 概念和价值。
                \begin{enumerate}
                  \item 定性。
                        \begin{enumerate}
                          \item \bs{方向决定道路，道路决定成败}。
                          \item 中法路是建法系，建法国的\bs{唯一正确道路}。
                        \end{enumerate}
                  \item 优势。
                        \begin{enumerate}
                          \item 中法路是\bs{最适合中国国情}的法治道路，
                          \item 根植于我国社会主义初级阶段的\zx{基本国情}，生发于我国\zx{改放}和\zx{社会主义现代化建设}的具体实践，
                          \item 是\bs{被实践证明了}的\zx{符合}我国基本\zx{国情}、符合人民\zx{群众愿望}、符合\zx{实践发展}要求的法治道路，
                          \item \bs{具有显著优越性}。
                        \end{enumerate}
                  \item 要求。
                        \begin{enumerate}
                          \item 要\zd{保持}在坚持中法路这个根本问题上的\zx{自信心和定力}，
                          \item 从中国实际出发，\zx{突出}中国特色、实践特色、时代特色，
                          \item \zx{既不能}罔顾国情、超越阶段，\zx{也不能}因循守旧、墨守成规，
                          \item 努力\zd{推进}各项法治工作的\zx{稳步向前}，借鉴但决\zx{不照搬}西方国家的法治经验和道路。
                        \end{enumerate}
                \end{enumerate}
          \item 如何（即中国特色社会主义法治的核心要义）
                \begin{enumerate}
                  \item 坚持中国共产\bs{党的领导}。略。
                  \item 坚持\bs{中特制}。
                        \begin{enumerate}
                          \item \blh 是中法系的\bs{根本制度基础}，是全推国的\bs{根本制度保障}。
                          \item \blh 是中国共产党领导人民在不断\zx{探索和实践}的基础上\zx{形成和发展}起来的。
                          \item 中特社根本制度、基本制度和重要制度，是中法路的\bs{制度基础和重要保障}。
                          \item 要坚持中法路，不断\zx{巩固和完善}中特制，\zx{以法治为}中特制\zx{保驾护航}。
                        \end{enumerate}
                  \item 坚持\bs{贯彻中法论}。
                        略（主要强调了以法治思想为代表的中国特色社会主义法治理论的重要性，法治思想的价值是昨天的内容）。
                \end{enumerate}
        \end{enumerate}
\end{enumerate}